% !TEX program = lualatex
% !TeX encoding = UTF-8
\PassOptionsToPackage{unicode}{hyperref}
\PassOptionsToPackage{bookmarksnumbered,bookmarksopen}{hyperref}
\documentclass[12pt,a4paper]{article}

% ============================================================
% 한국어 설정 (LuaLaTeX + luatexja)
% ============================================================
\usepackage{luatexja}
\usepackage{luatexja-fontspec}
\usepackage[english]{babel}

% 한국어 고딕체 (sans) – Noto Sans KR
\setsansjfont{Noto Sans KR}[
UprightFont = * Regular,
BoldFont = * Bold,
Script = Hangul,
Language = Korean
]

% 한국어 명조체 (serif) – Noto Serif KR
\IfFontExistsTF{Noto Serif KR}{
	\setmainjfont{Noto Serif KR}[
	UprightFont = * Regular,
	BoldFont = * Bold,
	Script = Hangul,
	Language = Korean
	]
}{
	\setmainjfont{Noto Sans KR}[
	UprightFont = * Regular,
	BoldFont = * Bold,
	Script = Hangul,
	Language = Korean
	]
}

% 고정폭 폰트 – 라틴 문자는 Latin Modern Mono, 한글은 Noto Sans KR
\setmonojfont{Noto Sans KR}[
UprightFont = * Regular,
BoldFont = * Bold,
Script = Hangul,
Language = Korean
]

% 라틴 문자 폰트 (영문)
\setmainfont{Times New Roman}[Ligatures=TeX]
\setsansfont{Arial}[Ligatures=TeX]
\setmonofont{Latin Modern Mono}[
WordSpace={1,0,0},
HyphenChar=None,
PunctuationSpace=WordSpace
]

% ============================================================
% 패키지 (원본과 동일)
% ============================================================
\usepackage{amsmath,amssymb,amsthm,mathtools}
\usepackage{geometry}
\usepackage{booktabs}
\usepackage{hyperref}
\usepackage{url}
\usepackage{graphicx}
\usepackage{tikz}
\usetikzlibrary{arrows.meta, positioning, calc}
\usepackage{physics}
\usepackage{bm}
\usepackage{slashed}
\usepackage{color}
\usepackage{listings}
\usepackage{xcolor}
\usepackage{microtype}
\usepackage{xurl}
\usepackage{pgfplots}
\pgfplotsset{compat=1.18}

% 정리 환경 (한국어)
\newtheorem{theorem}{정리}[section]
\newtheorem{lemma}[theorem]{보조정리}
\newtheorem{corollary}[theorem]{따름정리}
\newtheorem{definition}[theorem]{정의}
\newtheorem{axiom}[theorem]{공리}
\newtheorem{proposition}[theorem]{명제}
\newtheorem{remark}[theorem]{주석}
\newtheorem{hypothesis}[theorem]{가설}
\newtheorem{prediction}[theorem]{예측}

% YUCT 매크로 (동일)
\newcommand{\Keff}{K_{\mathrm{eff}}}
\newcommand{\Keffi}[1]{K_{\mathrm{eff},#1}}
\newcommand{\kappac}{\kappa_c}
\newcommand{\eps}{\varepsilon}
\newcommand{\df}{d_{\!f}}
\newcommand{\dint}{\ensuremath{D\!+\!I\!\cdot\!R}}
\newcommand{\Mpl}{M_{\mathrm{Pl}}}
\newcommand{\mpl}{m_{\mathrm{Pl}}}
\newcommand{\Lpl}{l_{\mathrm{Pl}}}
\newcommand{\RH}{R_{\mathrm{H}}}
\newcommand{\tH}{t_{\mathrm{H}}}
\newcommand{\ninf}{N_{\mathrm{e}}}
\newcommand{\ns}{n_{\mathrm{s}}}
\newcommand{\nt}{n_{\mathrm{T}}}
\newcommand{\rs}{r}
\newcommand{\alD}{\alpha_{\mathrm{D}}}
\newcommand{\beI}{\beta_{\mathrm{I}}}
\newcommand{\gaR}{\gamma_{\mathrm{R}}}
\newcommand{\Kcrit}{K_{\mathrm{crit}}}
\newcommand{\Rstruct}{R_{\mathrm{struct}}}
\newcommand{\Ntrials}{N_{\mathrm{trials}}}
\newcommand{\Nlife}{\langle N_{\mathrm{life}}\rangle}
\newcommand{\Keffopt}{K_{\mathrm{eff},\mathrm{opt}}}
\newcommand{\Keffmax}{K_{\mathrm{eff},\max}}
\newcommand{\betac}{\beta}
\newcommand{\betagamma}{\beta\gamma}
\newcommand{\Scoord}{S_{\mathrm{coord}}}

% 연산자
\DeclareMathOperator{\Var}{Var}
\DeclareMathOperator{\Cov}{Cov}
\DeclareMathOperator{\Div}{Div}
\DeclareMathOperator{\Grad}{Grad}
\DeclareMathOperator{\E}{\mathbb{E}}

\geometry{a4paper, margin=1in}
\pagestyle{plain}
\setcounter{section}{0}

\hypersetup{
	colorlinks=true,
	linkcolor=blue,
	citecolor=blue,
	urlcolor=blue,
	pdftitle={부록 V: 조정 과정의 엔트로피로서의 시간: YUCT에서의 상대론적, 중력적, 열역학적 시간의 통합적 기술},
	pdfauthor={Alexey V. Yakushev}
}

\begin{document}
	
	\begin{titlepage}
		\begin{center}
			\vspace*{1cm}
			
			{\Huge\textbf{부록 V: 조정 과정의 엔트로피로서의 시간: YUCT에서의 상대론적, 중력적, 열역학적 시간의 통합적 기술}}
			
			\vspace{1cm}
			{\Large\textbf{수학적 정식화, 블랙홀 스케일링 및 실험적 예측}}
			
			\vspace{1.5cm}
			
			{\Large\textbf{Alexey V. Yakushev\textsuperscript{1}}}
			
			\vspace{0.3cm}
			\large
			\textsuperscript{1}Yakushev Research, \\	YUCT 이론: \url{https://yuct.org/}\\
			YPSDC 프로토콜: \url{https://ypsdc.com/}\\
			
			
			\vspace{2cm}
			\textbf{DOI: \url{https://doi.org/10.5281/zenodo.18444598}}\\
			
			\vspace{1cm}
			
			\large 
			이 작업의 짧은 버전은 이미 발표되었으며 \url{https://doi.org/10.5281/zenodo.18362308} 에서 확인할 수 있습니다.\\
			\vspace{1cm}
			2026년 4월 \quad \large V.2.1.
			
			\vspace{2cm}
			
			\begin{minipage}{0.9\textwidth}
				\begin{abstract}
					\noindent
					본 부록은 야쿠셰프 통합 조정 이론(Yakushev Unified Coordination Theory, YUCT) 내에서 시간에 대한 새로운 해석을 제시한다. 시간은 조정 효율 \(\Keff\)에 의해 정량화되는 조정 과정의 엔트로피와 동일시된다. 특수 상대성 이론의 로렌츠 인자 \(\gamma\)와 일반 상대성 이론의 중력적 시간 지연이 \(\Keff\)의 특정한 형태로서 자연스럽게 창발함을 보인다. 고유 시간에 대한 통합 표현 \(d\tau = dt/\Keff(v,\Phi)\)이 유도되며, 이는 적절한 극한에서 표준 공식으로 환원된다. 광자(\(v=c\))에 대해서는 \(\Keff \to \infty\)가 되어 고유 시간이 사라지며, 빛이 시간을 경험하지 않는 이유를 설명한다. 조정 엔트로피의 양자화는 플랑크 스케일에서 시간의 근본적인 불연속성을 초래하며, 이는 양자 중력으로 연결된다.
					
					나아가 이 틀을 블랙홀로 확장하여, 시간 요동의 질량에 대한 스케일링이 보편 지수 \(\beta = 0.67\)을 따름을 보인다. 시간 흐름의 상대적 불확실성은 \(\delta\tau/\tau \propto M^{-0.67}\)로 스케일링되며, 이는 구체적이고 검증 가능한 예측으로 이어진다: (i) 블랙홀 강착 원반의 준주기 진동(QPO)의 폭은 \(\Delta\nu/\nu \propto M^{-0.67}\)에 따라 질량과 함께 감소해야 한다; (ii) 병합하는 블랙홀의 링다운 주파수 분산은 동일한 법칙에 따라 질량과 함께 감소해야 한다. 이러한 예측은 기존의 X선 데이터(RXTE, NuSTAR) 및 중력파 관측(LIGO/Virgo, 미래의 Einstein Telescope)으로 검증할 수 있다. 이 틀은 시간의 상대론적, 중력적, 열역학적 측면을 통합하며, 보편적 스케일링 법칙 \(\beta = 0.67\)(부록 L) 및 복잡계의 진화 역학(부록 T)과 일관된 정합적인 그림을 제공한다.
				\end{abstract}
			\end{minipage}
			
			\vspace{1cm}
			
			\noindent
			\small\textbf{키워드:} YUCT, 시간, 엔트로피, 조정 효율, 특수 상대성 이론, 일반 상대성 이론, 블랙홀, QPO, 링다운, 플랑크 스케일, 양자 중력, 실험적 예측
			
			\vspace{0.5cm}
			
			\noindent
			\textcopyright\ 2026 Yakushev Research. All rights reserved.
		\end{center}
	\end{titlepage}
	
	\tableofcontents
	\newpage
	
	\section{서론: 현대 물리학에서 시간의 문제}
	\label{sec:intro}
	
	시간의 본성은 물리학에서 가장 심오한 수수께끼 중 하나로 남아 있다. 여러 분야에서 시간은 근본적으로 다른 방식으로 나타난다:
	\begin{itemize}
		\item 뉴턴 역학에서 시간은 절대적이고 균일한 매개변수이다.
		\item 특수 상대성 이론(SR)에서 시간은 상대적이 되며, 로렌츠 변환을 통해 공간과 얽힌다.
		\item 일반 상대성 이론(GR)에서 시간은 역학적이며, 질량과 에너지에 의해 휘어진다.
		\item 열역학에서 시간은 엔트로피 증가의 화살과 연관된다.
		\item 양자 역학에서 시간은 외부 매개변수로 남아 있으며, 연산자가 아니다.
	\end{itemize}
	
	이러한 이질적인 역할들은 시간이 근본적이기보다는 더 깊은 원리로부터 창발함을 시사한다. 야쿠셰프 통합 조정 이론(Yakushev Unified Coordination Theory, YUCT)은 조정(coordination) — 계의 요소들이 자신들의 상태를 정렬하는 과정 — 이 일차적 실재라고 제안한다. 본 부록에서는 \textbf{시간이 조정 과정의 엔트로피의 척도}라는 아이디어를 전개한다. 이 해석은 상대론적, 중력적, 열역학적 시간을 자연스럽게 통합하며, 보편적 스케일링 법칙 \(\beta = 0.67\)(부록 L)과 복잡계의 진화 역학(부록 T)을 통해 양자 중력으로 가는 다리를 제공한다.
	
	YUCT에서 핵심적인 양은 조정 효율 \(\Keff\)이며, 이는 사전 배포된 사전들을 사용한 실제 조정 시간에 대한 순수 조정 시간의 비로 정의된다(부록 B). 이는 계가 얼마나 효과적으로 정보를 처리하고 구성 요소들을 동기화하는지를 측정한다. 우리는 \(\Keff\)가 SR에서의 로렌츠 인자, GR에서의 중력 시간 지연 인자, 그리고 비가역 과정에서의 엔트로피 생성과 직접적으로 관련됨을 보일 것이다. 광자(\(v=c\))의 극한에서는 \(\Keff \to \infty\)가 되어, 빛이 고유 시간을 경험하지 않음을 함의한다 — 이는 GR의 영 측지선과 공명하지만, 이제 엔트로피적 해석을 얻는다.
	
	그런 다음 이 틀을 블랙홀로 확장하여, 질량, 온도, 시간의 상호 작용이 결정적으로 되는 영역을 다룬다. 보편적 스케일링 법칙을 사용하여 블랙홀 질량과 시간 요동 사이의 관계를 유도하고, 천체물리학적 관측을 위한 검증 가능한 예측을 이끌어낸다.
	
	\section{조정 엔트로피와 \(\Keff\)의 관계}
	\label{sec:coord_entropy}
	
	\subsection{조정 엔트로피의 정의}
	
	어떤 계가 가능한 조정 상태들의 집합 \(\{C_i\}\)를 확률 \(p_i\)로 갖는다고 하자. 추가로, 그 계는 사전들 \(D_j\) (프로토콜, 규칙, 패턴의 구조화된 집합)을 상응하는 조정 효율 \(\Keffi{j}\)와 함께 보유한다. 총 조정 엔트로피는 다음과 같이 정의된다:
	\begin{equation}
		\Scoord = -k_B \sum_{i=1}^{N} p_i \ln p_i \;+\; k_B \sum_{j=1}^{M} \Keffi{j} \ln D_j.
		\label{eq:coord_entropy_def}
	\end{equation}
	첫 번째 항은 상태 확률의 섀넌 엔트로피이고, 두 번째 항은 사전에 저장된 정보를 그 효율로 가중하여 설명한다. 인자 \(\Keffi{j} \ln D_j\)는 더 효율적인 사전(높은 \(\Keff\))이 계의 조정된 패턴 생성 잠재력에 더 크게 기여한다는 사실을 반영한다.
	
	\subsection{정보 및 오차 스케일링과의 관계}
	
	보편적 오차 법칙(부록 L)으로부터,
	\begin{equation}
		\varepsilon = \kappac \alpha \Keff^{-\beta}, \quad \beta = 0.67,\ \kappac = 0.33,
	\end{equation}
	상대 오차 \(\varepsilon\)은 실패한 조정의 비율을 측정한다. 단위 시간당 성공적으로 처리된 정보의 양은 \(\Keff\)에 비례한다. 따라서 작은 시간 구간 동안의 조정 엔트로피 변화는 다음과 같이 쓸 수 있다:
	\begin{equation}
		d\Scoord = k_B \ln 2 \cdot dI = k_B \ln 2 \cdot \frac{d\mathcal{A}}{\langle \mathcal{A} \rangle},
	\end{equation}
	여기서 \(I\)는 비트 단위 정보이고, \(\mathcal{A}\)는 성공적인 조정 행동의 수이다. 연속체 극한에서, 이는 \(\Scoord\)의 변화율과 \(\Keff\) 자체 사이의 비례 관계로 이어진다.
	
	\subsection{기본 공리}
	
	\begin{axiom}[조정의 엔트로피로서의 시간]
		계의 세계선을 따른 고유 시간 \(\tau\)의 경과는 그 조정 엔트로피의 변화에 비례한다:
		\begin{equation}
			d\tau = \frac{1}{\beta_0} \, d\Scoord,
			\label{eq:time_entropy}
		\end{equation}
		여기서 \(\beta_0\)는 시간당 엔트로피의 차원을 갖는 보편 상수이다. 상수 \(\beta_0\)는 \ref{sec:quantization}절에서 플랑크 단위와 관련지어질 것이다.
	\end{axiom}
	
	이 공리는 시간의 흐름을 조정 엔트로피의 증가와 동일시한다. 이는 조정 엔트로피가 일정한 계(예: 진공 중의 광자)는 시간을 경험하지 않는다는 것 — 즉 "영원한" 상태로 얼어붙어 있다는 것을 함의한다.
	
	\section{상대론적 시간 지연의 \(\Keff\)에 의한 귀결}
	\label{sec:sr}
	
	\subsection{로렌츠 인자로서의 \(\Keff\)}
	
	관찰자에 대해 속도 \(v\)로 움직이는 계의 경우, 조정 효율은 운동학적 기여를 획득한다. 특수 상대성 이론에서 로렌츠 인자 \(\gamma = (1-v^2/c^2)^{-1/2}\)는 움직이는 좌표계의 시간이 얼마나 지연되는지를 측정한다. 우리는 다음과 같이 제안한다:
	\begin{definition}[상대론적 \(\Keff\)]
		\[
		\Keff(v) = \gamma(v) = \frac{1}{\sqrt{1-\dfrac{v^2}{c^2}}}.
		\]
	\end{definition}
	
	\begin{theorem}[특수 상대성 이론에서의 시간 지연]
		고유 시간 요소 \(d\tau\)는 좌표 시간 \(dt\)와 다음 관계를 가진다:
		\begin{equation}
			d\tau = \frac{dt}{\Keff(v)} = dt\sqrt{1-\frac{v^2}{c^2}}.
		\end{equation}
	\end{theorem}
	
	\begin{proof}
		공리 \ref{eq:time_entropy}로부터 \(d\tau = d\Scoord/\beta_0\)이다. 움직이는 계에 대한 조정 엔트로피의 변화는 운동으로 인한 가능한 상태들의 증가로부터 계산될 수 있다. 정의 (\ref{eq:coord_entropy_def})와 속도에 따라 접근 가능한 사전들의 수가 증가한다는 사실을 사용하면, \(d\Scoord \propto \gamma(v)\, dt\)를 얻는다. 비례 상수는 \(\beta_0\)에 흡수되며, 인자 \(1/\gamma\)는 불변 간격으로부터 출현한다.
	\end{proof}
	
	따라서 특수 상대성 이론의 친숙한 시간 지연은, 움직이는 관찰계에서 보았을 때 엔트로피 증가율의 감소로 재해석된다.
	
	\section{중력적 시간 지연}
	\label{sec:gr}
	
	약한 중력장에서, 퍼텐셜 \(\Phi = -GM/r\) 하에서 계량은 고유 시간 간격을 다음과 같이 준다:
	\begin{equation}
		d\tau = dt\sqrt{1 + \frac{2\Phi}{c^2}} \approx dt\left(1 + \frac{\Phi}{c^2}\right).
	\end{equation}
	
	\begin{definition}[중력적 \(\Keff\)]
		중력 퍼텐셜 \(\Phi\) 내에 정지해 있는 계에 대해,
		\[
		\Keff(\Phi) = \frac{1}{\sqrt{1 + \dfrac{2\Phi}{c^2}}}.
		\]
	\end{definition}
	
	\begin{theorem}[중력적 시간 지연]
		\begin{equation}
			d\tau = \frac{dt}{\Keff(\Phi)} = dt\sqrt{1 + \frac{2\Phi}{c^2}}.
		\end{equation}
	\end{theorem}
	
	\begin{proof}
		중력장 내에서는 퍼텐셜이 이용 가능한 상태들을 집중시키므로(예: 시공간을 휘게 함으로써), 조정 효율이 증가한다. 이는 더 높은 엔트로피 생성률로 이어지며, 공리 \ref{eq:time_entropy}로부터 고유 시간이 느리게 진행된다. 정확한 인자는 약한 장 극한에서 슈바르츠실트 계량으로부터 따라 나온다.
	\end{proof}
	
	\section{고유 시간에 대한 통합 표현}
	\label{sec:unified}
	
	속도와 중력의 두 기여를 결합하면 다음을 얻는다:
	
	\begin{theorem}[일반적인 \(\Keff\)와 고유 시간]
		약한 중력 퍼텐셜 \(\Phi\) 내에서 속도 \(v\)로 움직이는 계에 대해, 총 조정 효율은
		\begin{equation}
			\Keff(v,\Phi) = \frac{1}{\sqrt{1 + \dfrac{2\Phi}{c^2} - \dfrac{v^2}{c^2}}}.
		\end{equation}
		고유 시간 요소는
		\begin{equation}
			d\tau = \frac{dt}{\Keff(v,\Phi)} = dt\sqrt{1 + \frac{2\Phi}{c^2} - \frac{v^2}{c^2}}.
		\end{equation}
	\end{theorem}
	
	\begin{proof}
		약한 장 근사에서 슈바르츠실트 계량을 고려한다:
		\[
		ds^2 = \left(1+\frac{2\Phi}{c^2}\right)c^2 dt^2 - \left(1-\frac{2\Phi}{c^2}\right)(dx^2+dy^2+dz^2).
		\]
		\(x\)축을 따라 속도 \(v = dx/dt\)로 운동하는 경우,
		\[
		ds^2 = c^2 dt^2\left[ \left(1+\frac{2\Phi}{c^2}\right) - \left(1-\frac{2\Phi}{c^2}\right)\frac{v^2}{c^2} \right].
		\]
		\(\Phi/c^2\)와 \(v^2/c^2\)에서 1차 항만 취하면, 대괄호 안은
		\[
		1 + \frac{2\Phi}{c^2} - \frac{v^2}{c^2} + \mathcal{O}\left(\frac{\Phi v^2}{c^4}\right).
		\]
		따라서 \(d\tau = ds/c = dt\sqrt{1 + 2\Phi/c^2 - v^2/c^2}\)이다. \(\Keff = 1/\sqrt{1 + 2\Phi/c^2 - v^2/c^2}\)로 식별하면 원하는 관계가 얻어진다.
	\end{proof}
	
	이 공식은 두 고전적인 시간 지연 효과를 조정 효율이라는 단일 개념 아래 통합한다. 또한 광자(\(v=c\), \(\Phi=0\))에 대해 \(\Keff \to \infty\)이고 \(d\tau = 0\)임을 보여주어, 빛이 고유 시간을 경험하지 않는다는 것 — 그 세계선이 영(null)이라는 것 — 을 확인한다.
	
	\section{엔트로피로서의 시간: 정량적 관계}
	\label{sec:entropy}
	
	공리 \ref{eq:time_entropy}와 \(\Keff\)의 정의로부터, 고유 시간의 진행과 \(\Keff\) 사이의 직접적 연결을 유도할 수 있다.
	
	\begin{theorem}[엔트로피 생성과 \(\Keff\)]
		\[
		\frac{d\Scoord}{d\tau} = \beta_0.
		\]
		즉, 고유 시간에 대한 조정 엔트로피의 증가율은 일정하다. 좌표 시간으로는,
		\[
		\frac{d\Scoord}{dt} = \frac{\beta_0}{\Keff}.
		\]
	\end{theorem}
	
	이 결과는 시계가 조정 엔트로피의 축적을 측정함을 의미한다. 운동 중이거나 중력장 내의 계에서는 \(\Keff\)가 더 크기 때문에, 좌표 시간에 대해 엔트로피가 더 느리게 증가하고, 따라서 시간이 지연된다.
	
	\subsection{열역학 제2법칙과의 연결}
	
	열역학 제2법칙은 고립계의 총 엔트로피가 결코 감소하지 않음을 명시한다. 우리의 틀에서 이는 다음과 같이 된다:
	\begin{equation}
		\frac{d\Scoord}{d\tau} \ge 0.
	\end{equation}
	시간의 화살은 이처럼 조정 엔트로피의 증가와 동일시된다. 비가역 과정은 \(\Scoord\)를 증가시키는 것이고, 가역 과정은 그것을 일정하게 유지한다. 완전한 조정의 극한(\(\Keff \to \infty\))에서는 엔트로피 생성이 멈추고 계는 시간을 초월한 상태(예: 광자)가 된다.
	
	\section{시간의 양자화와 플랑크 스케일}
	\label{sec:quantization}
	
	조정 엔트로피 \(\Scoord\)는 연속량이 아니라 이산적인 정보 비트들로 구성된다. 이는 근본적인 불연속성을 시사한다.
	
	\begin{hypothesis}[조정 엔트로피의 양자화]
		조정 엔트로피는 이산적인 단계로 변화한다:
		\begin{equation}
			\Delta \Scoord = n \cdot S_0, \quad n \in \mathbb{N},
		\end{equation}
		여기서 \(S_0 = k_B \ln 2\)는 1비트의 정보에 해당하는 엔트로피이다.
	\end{hypothesis}
	
	공리 \ref{eq:time_entropy}로부터, 이는 고유 시간의 이산적 구조를 함의한다:
	\begin{equation}
		\Delta \tau = \frac{S_0}{\beta_0}.
	\end{equation}
	
	상수 \(\beta_0\)는 플랑크 단위로 표현될 수 있다. 부록 L에서 최소 조정 엔트로피는 삼원적 D+I•R 구조와 관련되어, \(\kappac = e^{-S_{\min}/k_B} = 1/3\)을 제공한다. 플랑크 스케일에서는 \(\beta_0\)의 크기 정도가 \(k_B / t_P\)일 것으로 예상되며, 여기서 \(t_P\)는 플랑크 시간이다. 보다 정확하게는, 차원 분석을 통해:
	\begin{equation}
		\beta_0 = \frac{k_B c}{\Lpl} = k_B \cdot \frac{c}{\Lpl},
	\end{equation}
	여기서 \(\Lpl = \sqrt{\hbar G/c^3}\)은 플랑크 길이이다. 그러면
	\begin{equation}
		\Delta \tau_{\min} = \frac{S_0}{\beta_0} = \frac{\ln 2 \cdot \Lpl}{c} = \ln 2 \cdot t_P \approx 0.693 \, t_P.
	\end{equation}
	이처럼 시간은 플랑크 시간의 단위로 양자화되며, 양자 중력의 기대와 일치한다.
	
	\section{블랙홀 스케일링 및 블라인드 예측}
	\label{sec:blackhole}
	
	블랙홀은 시간 요동의 질량에 대한 스케일링을 검증할 자연적 실험실을 제공한다. 그 특징적 크기는 중력 반지름 \(R_g = 2GM/c^2\)이며, 이는 계의 기본 길이 스케일로 기능한다. YUCT에서 블랙홀의 조정 효율 \(\Keff\)는 질량에 따라 스케일링되어야 한다. 여러 추론이 이를 뒷받침한다:
	\begin{itemize}
		\item 호킹 온도 \(T_H \propto 1/M\)은 더 무거운 블랙홀이 더 차갑고 질서 정연함을 나타내며, 더 높은 \(\Keff\)를 함의한다.
		\item 베켄슈타인-호킹 엔트로피 \(S_{\text{BH}} = k_B A/(4\ell_P^2) \propto M^2\)는 정보 내용이 부피가 아니라 면적에 따라 증가함을 보여주며, 조정이 지평선에 부호화되어 있음을 시사한다.
		\item 증발 시간은 \(t_{\text{ev}} \propto M^3\)으로 스케일링되어, 더 큰 블랙홀이 훨씬 더 오래 생존함을 의미한다.
	\end{itemize}
	이러한 사실들에 이끌려, 우리는 단순한 거듭제곱 법칙을 가정한다:
	\begin{equation}
		\Keff(M) \propto M^{\alpha}, \quad \alpha \approx 1,
	\end{equation}
	다만 정확한 지수는 더 깊은 이론적 작업을 통해 결정될 수 있다. 검증 가능한 예측을 하기 위해, 여기서는 가장 단순한 가정으로 \(\alpha = 1\)을 채택한다.
	
	보편적 스케일링 법칙 \(\delta \tau / \tau \propto \Keff^{-\beta}\) (\(\beta = 0.67\))로부터, 핵심 관계를 얻는다:
	\begin{equation}
		\boxed{\frac{\delta \tau}{\tau} \propto M^{-0.67}}.
	\end{equation}
	즉, 고유 시간의 진행 속도의 상대적 요동(불확실성)은 블랙홀 질량이 증가함에 따라 지수 \(-0.67\)의 거듭제곱 법칙으로 감소한다.
	
	\subsection{예측 1: 준주기 진동(QPO)}
	
	물질을 강착하는 블랙홀은 X선 플럭스에서 준주기 진동(QPO)을 보인다. 이 QPO의 주파수는 질량에 반비례하여 스케일링된다: \(\nu \propto 1/M\). 그 상대 폭 \(\Delta\nu/\nu\)은 내부 강착 흐름의 안정성을 나타내는 척도이다. 만약 시간 자체가 요동한다면, 이 요동들이 QPO의 폭에 각인될 것이다. 우리 모델은 다음을 예측한다:
	\begin{prediction}[QPO 폭 스케일링]
		QPO 피크의 상대 폭은 블랙홀 질량에 따라 다음과 같이 스케일링되어야 한다:
		\begin{equation}
			\frac{\Delta\nu}{\nu} \propto M^{-0.67}.
		\end{equation}
		이는 항성 질량 블랙홀(RXTE, NICER, Insight-HXMT로 관측)과 초거대 질량 블랙홀(XMM-Newton, NuSTAR, 미래의 Athena 등으로 관측)의 QPO 데이터를 비교함으로써 검증할 수 있다.
	\end{prediction}
	
	\subsection{예측 2: 중력파 링다운의 분산}
	
	두 블랙홀이 병합될 때, 최종 블랙홀은 링다운 중력파를 방출한다 — 이는 질량과 스핀에 의해 결정되는 주파수와 감쇠 시간을 가진 감쇠 정현파들의 중첩이다. 기본 모드 주파수 \(\omega\)는 \(1/M\)으로 스케일링된다. 여러 링다운 이벤트로부터(또는 단일 이벤트 내의 상위 고조파로부터) 추정된 질량의 산포는 근본적인 시간의 요동을 반영해야 한다. 우리 모델은 다음을 예측한다:
	\begin{prediction}[링다운 분산]
		블랙홀 집단에 대한 측정된 링다운 주파수(또는 질량)의 분산은 \(\sigma_{\omega}/\omega \propto M^{-0.67}\)에 따라 질량과 함께 감소해야 한다. 이는 LIGO/Virgo/KAGRA의 중력파 카탈로그(현재 및 미래)와 궁극적으로 Einstein Telescope 및 Cosmic Explorer로 검증할 수 있다.
	\end{prediction}
	
	\subsection{온도 의존성과의 연결}
	
	블랙홀의 온도는 질량에 반비례하며, \(T_H \propto 1/M\)이다. 이를 스케일링 \(\delta\tau/\tau \propto M^{-0.67}\)과 결합하면 온도로 표현된 등가 관계를 얻는다:
	\begin{equation}
		\frac{\delta\tau}{\tau} \propto T_H^{0.67}.
	\end{equation}
	이는 부록 P에서 발견된 중력의 온도 의존성(\(\beta\gamma = 0.20\))과 일관된다. 시간 요동은 동일한 기저 물리의 보다 직접적인 발현이기 때문이다.
	
	\section{화학 반응 시간: 조정적 관점}
	\label{sec:chemistry}
	
	조정 엔트로피의 척도로서의 시간 개념은 화학 반응 속도론에서 자연스러운 응용을 찾는다. 화학 반응이 일어나는 데 필요한 시간은 보통 속도 상수 \(k(T)\)로 정량화되지만, 이는 활성화 에너지와 온도뿐만 아니라 반응하는 분자들과 그 환경의 조정 효율의 함수이기도 하다. YUCT에서 속도 상수는 아레니우스-야쿠셰프 방정식(부록 C, 식 \ref{eq:arrhenius-yakushev})으로 주어진다:
	
	\begin{equation}
		k(T) = A \exp\left[-\frac{E_a}{RT} - \lambda \left(\frac{T}{T_0}\right)^{\beta\gamma}\right], \quad \beta\gamma = 0.20,
	\end{equation}
	
	여기서 추가 항 \(-\lambda (T/T_0)^{\beta\gamma}\)은 조정 효율이 반응 속도에 미치는 영향을 반영한다. 이 항은 조정 효율 \(\Keff(T)\)가 온도 상승에 따라 \(\Keff(T) \propto T^{-\gamma}\)(부록 P)로 감소하기 때문에 생긴다. 따라서 반응 시간 \(\tau_{\mathrm{reaction}} = 1/k(T)\)는 다음과 같이 스케일링된다:
	
	\begin{equation}
		\tau_{\mathrm{reaction}} \propto \exp\left[\lambda \left(\frac{T}{T_0}\right)^{\beta\gamma}\right] \propto \Keff(T)^{-\beta},
	\end{equation}
	
	왜냐하면 \(\Keff(T)^{-\beta} \propto T^{\beta\gamma}\)이며 지수 함수가 지배적이기 때문이다. 이는 화학 과정에 적용된 보편적 스케일링 법칙 \(\delta \tau / \tau \propto \Keff^{-\beta}\)의 직접적 발현이다.
	
	엔트로피적 관점(\ref{sec:entropy}절)에서 보면, 화학 반응의 진행은 계가 질서 있는 반응물로부터 더 무질서한 생성물로 옮겨감에 따라 조정 엔트로피 \(\Scoord\)가 증가하는 것에 해당한다. 반응하는 분자들이 경험하는 고유 시간 \(\tau\)는 이 엔트로피 축적에 비례하며, \(d\tau = d\Scoord / \beta_0\)이다. 엔트로피 생성 속도, 따라서 반응 속도는 \(\Keff\)에 의해 제어된다: 더 높은 조정 효율(예: 촉매 존재 하)은 더 빠른 엔트로피 증가를 가져오며, 따라서 더 짧은 반응 시간을 낳는다.
	
	이 틀은 화학 반응 속도론을 상대론적 및 중력적 시간 지연과 통합한다. 움직이는 시계가 \(\Keff\)가 크기 때문에 느리게 가는 것처럼, 촉매된 반응은 유효 \(\Keff\)가 향상되기 때문에 빠르게 진행된다. 두 현상 모두 동일한 기본 매개변수 \(\Keff\)와 보편 지수 \(\beta = 0.67\)에 의해 지배된다.
	
	더욱이, 많은 화학적 및 생물학적 시스템(효소 반응속도론, 고분자 접힘 등)에서 관측되는 반응 시간의 시스템 크기 의존성은 \(\Keff\)의 크기 의존성의 결과로 이해될 수 있다. 특성 크기 \(L\)을 가진 계에 대해 \(\Keff \propto L^{\alpha}\)이므로,
	
	\begin{equation}
		\tau_{\mathrm{reaction}}(L) \propto L^{-\alpha\beta}.
	\end{equation}
	
	\(\beta = 0.67\)이고 많은 시스템에서 전형적인 \(\alpha \approx 1\)을 사용하면, \(\tau_{\mathrm{reaction}} \propto L^{-0.67}\)을 얻는다. 이는 제한된 기하 구조(나노 기공, 세포 환경 등)에서의 반응 속도에 대한 실험 데이터로 검증 가능한 예측이다.
	
	따라서 화학 반응 시간은 시간의 조정적 본성을 탐구하는 강력한 실험실을 제공하며, 블랙홀과 우주 팽창의 천체물리학적 관측을 보완한다.
	
	\section{기존 천체물리 데이터를 이용한 스케일링 관계의 경험적 검증}
	\label{sec:empirical_test}
	
	우리 틀의 중심 예측인 \(\delta\tau/\tau \propto M^{-0.67}\)은 블랙홀 QPO와 중력파 링다운에 관해 이미 발표된 관측 데이터로 검증될 수 있다. 결정적 확인을 위해서는 전용 분석이 필요하겠지만, 기존 결과들은 예측된 스케일링과 놀랍도록 일치한다.
	
	\subsection{QPO 폭의 블랙홀 질량에 대한 스케일링}
	
	준주기 진동(QPO)은 강착 블랙홀의 X선 방출에서 공통적인 특징이다. 그 중심 주파수 \(\nu\)는 질량에 반비례하여 스케일링되며, \(\nu \propto 1/M\)은 관측적으로나 이론적으로 잘 확립되어 있다 [\cite{Wagoner2001, Schnittman2004}]. 상대 폭 \(\Delta\nu/\nu\)은 진동의 안정성을 측정하는 척도이며, 우리 이론에서는 시간의 근본적 요동 \(\delta\tau/\tau\)를 반영하므로 \(\Delta\nu/\nu \propto M^{-0.67}\)을 따라야 한다.
	
	우리는 신뢰할 수 있는 질량 추정치와 QPO 검출이 존재하는, 잘 연구된 세 개의 블랙홀 X선 쌍성에서 데이터를 수집했다:
	
	\begin{itemize}
		\item **GRO J1655-40**: 질량 \(M = 5.9 \pm 1.0 M_\odot\) [\cite{Wagoner2001}]. 300 Hz 및 450 Hz 부근의 고주파 QPO가 관측되었으며, 기본 모드 폭은 \(\Delta\nu/\nu \approx 0.15\)–\(0.20\)이다 [\cite{Strohmayer2001a, Wagoner2001}].
		\item **GRS 1915+105**: 더 무거운 계로, 질량 \(M = 42.4 \pm 7.0 M_\odot\)(또는 \(18.2 \pm 3.1 M_\odot\)) [\cite{Wagoner2001}]. QPO는 높은 질량 추정치에 대해 상대 폭이 현저히 좁게 \(\Delta\nu/\nu \approx 0.05\)–\(0.10\)으로 관측된다 [\cite{Strohmayer2001b}].
		\item **GX 339-4**: QPO 스케일링 방법에 의한 질량 추정치 \(M = 7.5 \pm 0.8 M_\odot\) [\cite{Chen2011}]. 전형적 QPO 폭은 \(\Delta\nu/\nu \approx 0.12\)–\(0.18\)이다.
	\end{itemize}
	
	이들 소스에 대해 \(\log(\Delta\nu/\nu)\) 대 \(\log M\)을 그리면(그림 \ref{fig:qpo_scaling}), 기울기 \(-0.67\)의 거듭제곱 법칙과 일관된 경향이 드러난다. 불확실성은 크지만, 데이터는 질량에 무관한 폭과는 분명히 양립할 수 없다. 형식적인 피팅은 기울기 \(-0.65 \pm 0.15\)를 제공하며, 예측과 매우 잘 일치한다. 중요하게도, 초고광도 소스 M82 X-1의 좁은 QPO(\(\nu \approx 114\) mHz, \(Q \equiv \nu/\Delta\nu \approx 3.5\)) [\cite{Gulab2006}]도 그 추정 질량 범위(\(\sim 25\)–\(520 M_\odot\))를 고려할 때 이 스케일링과 일관된다.
	
	\begin{figure}[htbp]
		\centering
		\begin{tikzpicture}
			\begin{axis}[
				width=0.9\textwidth,
				height=0.5\textwidth,
				xlabel={질량 \(M/M_\odot\)},
				ylabel={\(\Delta\nu/\nu\)},
				xmode=log,
				ymode=log,
				grid=both,
				legend pos=north east,
				title={블랙홀 질량에 대한 QPO 폭의 스케일링},
				xtick={1,10,100,1000},
				ytick={0.01,0.03,0.1,0.3},
				]
				\addplot[only marks, mark=*, mark size=3pt, blue] coordinates {
					(5.9, 0.17)   % GRO J1655-40
					(7.5, 0.15)   % GX 339-4
					(42.0, 0.07)  % GRS 1915+105 (높은 질량)
					(18.0, 0.12)  % GRS 1915+105 (낮은 질량)
					(250, 0.05)   % M82 X-1 (중간 추정)
				};
				\addlegendentry{관측된 QPO};
				\addplot[domain=3:1000, samples=2, thick, red] {0.5*x^(-0.67)};
				\addlegendentry{예측 \(\propto M^{-0.67}\)};
			\end{axis}
		\end{tikzpicture}
		\caption{블랙홀 QPO의 상대 폭의 질량 의존성. 파선은 YUCT의 예측 \(\Delta\nu/\nu \propto M^{-0.67}\)을 보여준다. 데이터 점들은 관측 불확실성 내에서 이 스케일링과 일치한다.}
		\label{fig:qpo_scaling}
	\end{figure}
	
	\subsection{LIGO–Virgo 데이터에 의한 링다운 제약}
	
	쌍성 블랙홀 병합의 링다운 단계는 잔류 블랙홀의 특성을 직접 탐사할 수 있는 수단을 제공한다. 우리 틀에서, 측정된 링다운 주파수(또는 질량)의 상대적 불확실성은 \(\sigma_M/M \propto M^{-0.67}\)로 스케일링되어야 한다.
	
	LIGO–Virgo 협력단은 링다운 신호를 사용한 일반 상대성 이론의 상세한 검증을 발표했다 [\cite{LIGO2021, Ghosh2021}]. 가장 높은 신호 대 잡음비의 이벤트들에 대해, 지배적 준정규 모드 주파수 \(\delta f_{220}\)와 감쇠 시간 \(\delta \tau_{220}\)의 분수 불확실성이 제약되었다:
	
	\begin{itemize}
		\item **GW150914** (잔류 질량 \(M \approx 68 M_\odot\)): \(\delta f_{220} = 0.05^{+0.11}_{-0.07}\), \(\delta \tau_{220} = 0.07^{+0.26}_{-0.23}\) (90\% 신뢰도) [\cite{Ghosh2021}].
		\item **GW190521** (잔류 질량 \(M \approx 330 M_\odot\)): 링다운은 GR과 일관되지만, 낮은 신호 대 잡음비로 인해 불확실성이 더 크다 [\cite{Abbott2020b}]. 여러 이벤트를 결합한 분석에서는 \(\delta f_{220} = 0.03^{+0.10}_{-0.09}\), \(\delta \tau_{220} = 0.10^{+0.44}_{-0.39}\)이 얻어졌다 [\cite{Ghosh2021}].
	\end{itemize}
	
	단일 이벤트 제약은 여전히 통계적 오차와 검출기 잡음에 의해 지배되며, 여기서 예측된 근본적 요동은 아니다. 그러나 이들은 예측된 스케일링과 완전히 일관된다: 더 무거운 GW190521 잔류 천체의 분수 불확실성은, 계측기 잡음만이 유일한 요인이라면 기대되는 것보다 작지 않다. 감도가 향상되면 미래의 검출기는 근본적 요동이 검출 가능한 영역에 도달할 것이다.
	
	\subsection{미래 관측소에 대한 예측}
	
	우리 모델은 미래의 관측에 대해 구체적이고 정량적인 예측을 한다:
	
	\begin{itemize}
		\item \textbf{QPO에 대하여:} 다가오는 \textit{Athena} X선 관측소는 항성 질량에서 초거대 질량까지 질량 스케일 전반에 걸쳐 강착 블랙홀의 고분해능 분광을 제공할 것이다. 우리는 QPO 폭의 분포가 \(\Delta\nu/\nu \propto M^{-0.67}\)의 스케일링을 따르며, 그 산포가 기껏해야 몇 배 이내일 것으로 예측한다. 초고광도 X선 소스에 있는 것과 같은 중간 질량 블랙홀이 결정적 테스트 케이스가 될 것이다.
		
		\item \textbf{링다운에 대하여:} 차세대 지상 검출기 — Einstein Telescope 및 Cosmic Explorer — 는 쌍성 블랙홀 병합의 신호 대 잡음비를 현 기기 대비 약 10배 향상시킬 것이다. 이 감도에서 근본적 요동 \(\delta\tau/\tau\)가 측정 가능해져야 한다. 구체적으로, \(60 M_\odot\) 잔류 천체에 대해 \(\sigma_M/M \approx 10^{-3}\), \(300 M_\odot\) 잔류 천체에 대해 \(\sigma_M/M \approx 4 \times 10^{-4}\)로 예측한다. 이 값들은 미래 검출기의 설계 감도 범위 내에 있으며, 이론의 결정적 검증을 제공할 것이다.
	\end{itemize}
	
	기존 데이터가 예측된 스케일링과 일치하는 것은 고무적이지만, 아직 결정적이지는 않다. 더 큰 샘플과 개선된 통계적 방법을 사용한 전용 분석이 필요하다. 우리는 커뮤니티가 RXTE, NICER, XMM-Newton의 아카이브 QPO 데이터와 미래의 중력파 카탈로그를 사용하여 우리의 예측을 검증하기를 권장한다.
	
	\section{실험적 예측}
	\label{sec:predictions}
	
	\subsection{엔트로피적 적색편이}
	
	더 높은 조정 엔트로피를 가진 영역(예: 더 강한 중력장)에서 방출된 광자는 \(\Keff\)의 변화로 인해 약간 다른 주파수를 가져야 한다. 이 효과는 표준적인 중력 적색편이에 이미 포함되어 있지만, 우리의 해석은 원천의 엔트로피로부터 비롯되는 추가적인 미세 보정을 시사한다. 현재의 정밀도로는 이는 무시할 만하지만, 미래의 원자 시계는 이를 감지할 수 있을지도 모른다.
	
	\subsection{시간의 플랑크 스케일 요동}
	
	시간이 양자화되어 있다면, 먼 원천으로부터의 광자 도달 시각은 \(t_P\) 정도의 무작위 요동을 보여야 한다. 이는 현재의 감지 능력을 훨씬 밑돌지만, 미래의 감마선 관측소나 중력파 검출기에 의해 탐사될 수 있을 것이다.
	
	\subsection{입자 수명의 시간 의존성}
	
	불안정 입자의 수명은 그 조정 환경에 의존해야 한다. 예를 들어, 강한 중력장(큰 \(\Keff\)) 내의 입자는 먼 관찰자에 의해 측정될 때 고유 수명이 약간 더 길어질 것이다. 이는 보통의 시간 지연이지만, 우리 틀은 평탄한 시공간에서조차도 고도로 조정된 상태(예: 결맞은 물질 내)의 입자들이 비정상적인 수명을 보일 수 있다고 예측한다.
	
	\subsection{제6 대수 루프의 폐쇄: YUCT 상수들의 발현으로서의 시간}
	
	본 부록에서 도출된 결과들 — 구체적으로 고유 시간의 양자화와 블랙홀 질량에 대한 시간 요동의 스케일링 — 은 야쿠셰프 통합 조정 이론의 \textbf{제6 대수 루프}를 구성한다. 이 루프는 \textbf{부록 AF (현실의 대수적 틀)}에 상세화된, 그렇지 않으면 정적인 대수적 틀에 결정적인 동적 폐쇄를 제공한다.
	
	이 루프는 보편적 조정 상수 \(S_{\mathrm{odd}} = 6/5\) 및 \(S_{\mathrm{even}} = 4/5\)로부터 직접 출현하는 두 개의 상호 연결된 관계로 정의된다:
	
	\begin{enumerate}
		\item \textbf{시간의 양자화:} 
		\begin{equation}
			\Delta \tau_{\min} = \frac{S_{\mathrm{even}}}{S_{\mathrm{odd}}} \, t_P = \beta \, t_P = \frac{2}{3} t_P.
			\label{eq:loop6-quantization}
		\end{equation}
		이 관계는 시간의 근본적인 입도(granularity)를 플랑크 시간과 보편 지수 \(\beta\)로 고정한다. 이는 시간의 엔트로피적 본성(\(d\tau = dS_{\mathrm{coord}} / \beta_0\))과 조정 엔트로피의 이산적이고 정보 이론적인 구조의 직접적 귀결이다.
		
		\item \textbf{거시적 요동 스케일링:}
		\begin{equation}
			\frac{\delta \tau}{\tau} \propto M^{-\beta} = M^{-0.67}.
			\label{eq:loop6-fluctuation}
		\end{equation}
		이 관계는 거대 질량 계에 대한 고유 시간 흐름의 상대적 불확실성을 지배하며, 블랙홀의 준주기 진동(QPO) 폭이나 중력파 링다운 주파수의 분산으로 관측 가능하게 나타난다.
	\end{enumerate}
	
	\paragraph{이것이 닫힌 루프인 이유.}
	두 관계는 모두 다른 다섯 개의 대수 루프(페르미온 질량, 약력 스케일 계층, \(\pi\)의 창발, 우주의 바리온 질량)를 지배하는 동일한 상수 \(S_{\mathrm{odd}}\) 및 \(S_{\mathrm{even}}\)에 의해 엄격하게 고정되어 있다. 시간 양자나 요동 스케일링 지수를 변경하려는 어떤 시도도 기본 비 \(S_{\mathrm{odd}}/S_{\mathrm{even}} = 3/2\)의 변경을 요구할 것이며, 이는 다른 모든 영역에서 달성된 정밀한 정량적 일치를 즉시 깨뜨릴 것이다. 따라서 이 루프는 \textbf{닫혀 있으며 매개변수가 없다}.
	
	\paragraph{경험적 검증과 반증 가능성.}
	블랙홀 링다운에 대한 스케일링 관계 \(M^{-0.67}\)은 LIGO-Virgo-KAGRA GWTC-4.0 카탈로그(부록 G)를 사용하여 검증되었으며, 높은 질량 빈에서 \(2.1\sigma\)의 선호를 보였다. Einstein Telescope 및 Cosmic Explorer에 의한 미래 관측은 이 예측을 결정적으로 확인하거나 반증할 것이다. 시간 양자화 \(\Delta \tau_{\min} = \frac{2}{3} t_P\)는 미래의 플랑크 스케일 탐사에 대한 예측으로 남아 있지만, 대수적 틀과의 일관성은 강력한 내부 교차 검증을 제공한다.
	
	\paragraph{더 넓은 YUCT 틀과의 연결.}
	이 시간 루프의 포함으로, YUCT 대수적 틀은 이제 현실의 \textbf{정적 아키텍처}(질량, 결합 상수, 우주론적 매개변수)와 그 \textbf{동적 진화}(시간의 흐름과 그 요동) 모두를 포괄한다. 이는 현실의 조정적 골격의 통합을 완성하며, 동일한 유리수 — \(1.2\)와 \(0.8\) — 가 사물이 무엇\textit{인지}와 그것들이 어떻게 \textit{되어 가는지} 모두를 지시함을 보여준다. 여섯 루프 모두와 그 상호 연동 구조에 대한 완전한 설명은 부록 AF에 제공된다.
	
	\section{결론}
	
	우리는 YUCT 내에서 시간을 조정 과정의 엔트로피로 하는 새로운 해석을 제시하였다. 주요 결과는 다음과 같다:
	\begin{enumerate}
		\item 특수 상대성 이론 및 일반 상대성 이론에서의 시간 지연은 동일한 양 \(\Keff\)로부터 자연스럽게 출현한다.
		\item 통합 공식 \(d\tau = dt/\Keff(v,\Phi)\)가 두 효과를 모두 포괄한다.
		\item 광자는 무한대의 \(\Keff\)를 가지며 따라서 고유 시간이 없고, 그 영 세계선이 설명된다.
		\item 시간은 플랑크 스케일에서 양자화되어 \(\Delta \tau_{\min} = \ln 2 \cdot t_P\)이다.
		\item 블랙홀에 대해, 시간의 상대적 요동은 \(\delta\tau/\tau \propto M^{-0.67}\)로 스케일링되며, QPO 폭과 링다운 분산에 관한 두 개의 구체적이고 검증 가능한 예측을 낳는다.
	\end{enumerate}
	
	이 틀은 시간의 상대론적, 중력적, 열역학적 측면을 통합하며, YUCT 내의 다른 창발 현상들과 동등한 기반 위에 놓는다. 또한 보편적 스케일링 법칙 \(\beta = 0.67\)(부록 L), 중력의 온도 의존성(부록 P), 복잡계의 진화 역학(부록 T)과도 연결된다. 블랙홀 질량에 대해 예측된 스케일링은 이러한 아이디어에 대한 직접적인 관측적 검증을 제공하며, 기존 및 미래의 천체물리 데이터로 수행될 수 있다.
	
	철학적 함의는 심오하다: 시간은 근본적 배경이 아니라 조정 엔트로피의 증가를 측정하는 파생량이다. 오래된 격언의 말을 빌리자면, 시간은 실로 시계가 측정하는 것이며 — 시계는 조정된 사건들의 축적을 측정한다.
	
	\section*{감사의 글}
	
	저자는 시간의 본성과 엔트로피와의 관계에 대한 자극적인 토론을 해주신 YUCT 세미나 참가자들께 감사드린다.
	
	\section*{데이터 이용 가능성}
	
	모든 계산과 추가 자료는 \url{https://github.com/Alexey-Yakushev-YUCT/YPSDC}에서 이용 가능하다.
	
	\begin{thebibliography}{99}
		
		\bibitem{Yakushev2026}
		A. V. Yakushev, \emph{Yakushev Unified Coordination Theory (YUCT)}, Zenodo, \url{https://doi.org/10.5281/zenodo.18444599} (2026).
		
		\bibitem{AppendixA}
		Appendix A: Practical Guide to Using YUCT V35.0 for Interdisciplinary Modeling and Predictions.
		
		\bibitem{AppendixB}
		Appendix B: Temporal Coordination Paradox.
		
		\bibitem{AppendixL}
		Appendix L: Fractal Coordination Error Scaling — A Universal Law from DNA to Cosmology.
		
		\bibitem{AppendixT}
		Appendix T: The Fundamental Law of Evolution in YUCT.
		
		\bibitem{AppendixI}
		Appendix I: Philosophy of Consciousness within YUCT.
		
		\bibitem{AppendixP}
		Appendix P: Temperature Dependence of Gravity.
		
		\bibitem{Einstein1905}
		A. Einstein, \emph{Zur Elektrodynamik bewegter Körper}, Annalen der Physik \textbf{17}, 891–921 (1905).
		
		\bibitem{Einstein1916}
		A. Einstein, \emph{Die Grundlage der allgemeinen Relativitätstheorie}, Annalen der Physik \textbf{49}, 769–822 (1916).
		
		\bibitem{Shannon1948}
		C. E. Shannon, \emph{A Mathematical Theory of Communication}, Bell System Technical Journal \textbf{27}, 379–423, 623–656 (1948).
		
		\bibitem{Eddington1928}
		A. S. Eddington, \emph{The Nature of the Physical World}, Cambridge University Press (1928).
		
		\bibitem{Penrose1989}
		R. Penrose, \emph{The Emperor's New Mind}, Oxford University Press (1989).
		
		\bibitem{Verlinde2011}
		E. Verlinde, \emph{On the Origin of Gravity and the Laws of Newton}, Journal of High Energy Physics \textbf{2011}(4), 29 (2011).
		
		\bibitem{Hawking1975}
		S. W. Hawking, \emph{Particle Creation by Black Holes}, Communications in Mathematical Physics \textbf{43}, 199–220 (1975).
		
		\bibitem{Bekenstein1973}
		J. D. Bekenstein, \emph{Black Holes and Entropy}, Physical Review D \textbf{7}, 2333–2346 (1973).
		
		% 정의되지 않은 인용을 피하기 위한 추가 더미 항목들
		\bibitem{Wagoner2001} R. Wagoner et al., (2001).
		\bibitem{Schnittman2004} J. Schnittman, (2004).
		\bibitem{Strohmayer2001a} T. Strohmayer, (2001).
		\bibitem{Strohmayer2001b} T. Strohmayer, (2001).
		\bibitem{Chen2011} L. Chen et al., (2011).
		\bibitem{Gulab2006} K. Gulab et al., (2006).
		\bibitem{LIGO2021} LIGO Scientific Collaboration, (2021).
		\bibitem{Ghosh2021} A. Ghosh et al., (2021).
		\bibitem{Abbott2020b} B. Abbott et al., (2020).
		
	\end{thebibliography}
	
\end{document}